% Propietario conservado: /Users/ruben/Documents/New project/output/EXTREMA_MEDIA_RAZON_RECIPROCIDAD_ES_EN_20260918/ARTICULO/ES/source/sections/iv_elipse_radio.tex
% RESULTADO_RECUPERADO. II REV10, 07_elipse.tex:1--61,
% 07b_geometria_elipse.tex:73--115 y 12_dualidad_t.tex:27--107.
% Explicitaciones algebraicas: cota H5<2 para verificar la carta y
% sustitución Rstar^4=(499 alpha-eta_ret)/(501 alpha+eta_ret).
\section{Action ellipse and determination of the normalized radius}
\label{x_reciprocidad:iv:ellipse:section}

The angular reader and the action section preserve the same
APP--TRIT--TPK antecedent. The form constructed
below uses the publications
\eqref{x_reciprocidad:iv:angular:representantes} and a positive section
\(\hbar=\hbar_s\), with
\(s\in\{\mathrm{pre},\mathrm{ret}\}\), of
\eqref{x_reciprocidad:iv:action:hbar}. The angular ratio will determine the form;
the action section will preserve its common scale.

\subsection{Domain and labelled semiaxes}
\label{x_reciprocidad:iv:ellipse:dominio}

For brevity, let \(A=A_{\deg}\) and \(C^*=C^*_{\deg}\).
In the chart \(A\ne0\), \(\hbar>0\) and \(|C^*/A|<1\),
we define
\begin{equation}
 \xi=\frac{C^*}{A},\qquad
 \eta_{\mathrm{el}}=\operatorname{artanh}\xi,\qquad
 a_S=\sqrt{2\hbar}\,e^{\eta_{\mathrm{el}}/2},\qquad
 b_S=\sqrt{2\hbar}\,e^{-\eta_{\mathrm{el}}/2}.
 \label{x_reciprocidad:iv:ellipse:semiejes}
\end{equation}
The semiaxes and the balanced coordinates \(Q,P\) have
units of the square root of action. The axes remain labelled:
the sign of \(\eta_{\mathrm{el}}\) determines which is
larger. The curve and the product of its semiaxes are
\begin{equation}
 \gamma(\theta)=(a_S\cos\theta,b_S\sin\theta),\qquad
 \frac{Q^2}{a_S^2}+\frac{P^2}{b_S^2}=1,\qquad
 a_Sb_S=2\hbar.
 \label{x_reciprocidad:iv:ellipse:curva}
\end{equation}

\begin{proposition}[Membership of the publications in the elliptic chart]
\label{x_reciprocidad:iv:ellipse:pertenencia}
The positive branch of \eqref{x_reciprocidad:iv:angular:representantes}, with the
intervals \eqref{x_reciprocidad:iv:action:intervalos}, satisfies
\(A>0\), \(0<C^*<A\), and \(\hbar_s>0\).
Reversal of the odd coordinate preserves \(|C^*/A|<1\).
\end{proposition}
\begin{proof}
Proposition \ref{x_reciprocidad:iv:action:positividad} gives
\(0<\eta_{\mathrm{ret}}<H_5\) and positivity of both
action sections. For an upper bound on \(H_5\),
\eqref{x_reciprocidad:iv:action:intervalos} implies
\(\alpha<1/125\) and \(\varphi>8/5\). Retaining only
the positive terms of \eqref{x_reciprocidad:iv:action:h5} yields
\[
 H_5<
 \frac{\sqrt5}{2}+
 \frac{9}{16\cdot125^2}+
 \frac{7}{48\cdot125^4}.
\]
Since \(5<81/16\), one has \(\sqrt5/2<9/8\). Moreover,
\[
 \frac{9}{16\cdot125^2}+
 \frac{7}{48\cdot125^4}
 =\frac{421882}{11718750000}<\frac1{1000}.
\]
Thus \(H_5<9/8+1/1000<2\). Using
\(\alpha>7/1000\) now gives
\[
 0<C^*=2(\eta_{\mathrm{ret}}+\alpha)
 <2\left(2+\frac1{125}\right)=\frac{502}{125}
 <7<1000\alpha=A.
\]
The channel interchange in
\ref{x_reciprocidad:iv:angular:inversion} replaces \(C^*\) by \(-C^*\)
and preserves \(A\); it therefore preserves the absolute inequality.
\end{proof}

\subsection{Parametric phase and oriented area}
\label{x_reciprocidad:iv:ellipse:area}

\begin{proposition}[Area law and transport of angular coordinates]
\label{x_reciprocidad:iv:ellipse:ley-areal}
The parametric phase \(\theta\) of \(\gamma\) and its polar
angle \(\phi\), continuously lifted from zero, satisfy
\begin{equation}
 \phi(\theta)=\arg(a_S\cos\theta+i b_S\sin\theta),\qquad
 \frac{d\phi}{d\theta}
 =\frac{a_Sb_S}{a_S^2\cos^2\theta+b_S^2\sin^2\theta}>0.
 \label{x_reciprocidad:iv:ellipse:fases}
\end{equation}
In a chart where both tangents are defined,
\(\tan\phi=(b_S/a_S)\tan\theta\). The oriented area swept out
from phase zero is
\begin{equation}
 \mathcal A(\theta)=
 \frac12\int_0^\theta(QP'-PQ')\,dt=\hbar\theta,
 \qquad
 \mathcal A(1)=\hbar,\quad\mathcal A(2\pi)=h_s.
 \label{x_reciprocidad:iv:ellipse:area-fase}
\end{equation}
\end{proposition}
\begin{proof}
The derivative of the argument of a nonzero vector is
\((QP'-PQ')/(Q^2+P^2)\). For the curve
\eqref{x_reciprocidad:iv:ellipse:curva}, the numerator equals
\(a_Sb_S=2\hbar\); this proves the derivative formula and its
positivity. The quotient \(P/Q\) gives the tangent relation;
the continuous lift preserves quadrants and turns.
The integral equals \(a_Sb_S\theta/2=\hbar\theta\).
Finally, \(h_s=2\pi\hbar_s\) gives the identity for
the complete turn.
\end{proof}

One radian of parametric phase sweeps an area \(\hbar\). The polar
angle retains its own derivative in
\eqref{x_reciprocidad:iv:ellipse:fases}. In polar coordinates, the equivalent
law is
\[
 d\mathcal A=\frac12\rho_{\partial}(\phi)^2\,d\phi,\qquad
 \rho_{\partial}(\phi)=
 \left(\frac{\cos^2\phi}{a_S^2}
       +\frac{\sin^2\phi}{b_S^2}\right)^{-1/2}.
\]
The radial formula follows by substituting
\((Q,P)=\rho(\cos\phi,\sin\phi)\) into
\eqref{x_reciprocidad:iv:ellipse:curva}. With the orientation of
\(\gamma\), the integral of the canonical one-form retains
the sign
\begin{equation}
 \oint_\gamma P\,dQ=-\pi a_Sb_S=-h_s.
 \label{x_reciprocidad:iv:ellipse:accion-orientada}
\end{equation}
Indeed, \(P\,dQ=-a_Sb_S\sin^2\theta\,d\theta\), and
the integral of \(\sin^2\theta\) over one turn is \(\pi\).

\subsection{Shape reciprocity and symplectic orientation}
\label{x_reciprocidad:iv:ellipse:reciprocidad}

Define the matrices
\begin{equation}
 D(\eta)=\begin{pmatrix}e^\eta&0\\0&e^{-\eta}\end{pmatrix},
 \qquad
 S_2=\begin{pmatrix}0&1\\1&0\end{pmatrix},
 \qquad
 J_2=\begin{pmatrix}0&-1\\1&0\end{pmatrix}.
 \label{x_reciprocidad:iv:ellipse:transportes}
\end{equation}

\begin{proposition}[Quadratic form and action under reciprocity]
\label{x_reciprocidad:iv:ellipse:accion-reciproca}
Both transports satisfy
\(T^{\mathsf T}D(-\eta)T=D(\eta)\).
For \(\omega=dQ\wedge dP\), the interchange \(S_2\) is
antisymplectic and the rotation \(J_2\) is symplectic:
\[
 S_2^*\omega=-\omega,\qquad J_2^*\omega=\omega.
\]
If \(\mathcal I(\gamma)=\oint_\gamma P\,dQ\) on a
closed curve, then
\begin{equation}
 \mathcal I(S_2\gamma)=-\mathcal I(\gamma),\qquad
 \mathcal I(J_2\gamma)=\mathcal I(\gamma).
 \label{x_reciprocidad:iv:ellipse:signos-accion}
\end{equation}
\end{proposition}
\begin{proof}
Matrix multiplication interchanges the two diagonal entries
of \(D\), proving both congruences.
The determinants of \(S_2,J_2\) are \(-1,1\);
this gives the transformation of the area form.
For \(\lambda=P\,dQ\), direct calculation gives
\[
 S_2^*\lambda=d(PQ)-\lambda,\qquad
 J_2^*\lambda=\lambda-d(PQ).
\]
The integral of the exact differential over the closed curve
is zero, proving both action identities.
\end{proof}

Orientation distinguishes the two transports even though they express
the same quadratic form, including in the circular case.
The labelling of the axes and the sign of the action therefore belong
to the data preserved in passing to the
lattice representation.

\Needspace{12\baselineskip}
\subsection{Normalized radius and lattice-readout matrix}
\label{x_reciprocidad:iv:ellipse:radio}

\begin{proposition}[Determination and recoverability of the shape radius]
\label{x_reciprocidad:iv:ellipse:radio-prop}
The labelled ellipse determines a unique positive radius
\begin{equation}
 R_\star=\sqrt{\frac{b_S}{a_S}}
 =e^{-\eta_{\mathrm{el}}/2}
 =\left(\frac{A-C^*}{A+C^*}\right)^{1/4},
 \qquad \tau_{\mathrm{el}}=iR_\star^2.
 \label{x_reciprocidad:iv:ellipse:rstar}
\end{equation}
The angular ratio is recovered through
\begin{equation}
 \frac{C^*}{A}=\frac{1-R_\star^4}{1+R_\star^4},
 \qquad
 R_\star^4=
 \frac{499\alpha-\eta_{\mathrm{ret}}}
      {501\alpha+\eta_{\mathrm{ret}}}.
 \label{x_reciprocidad:iv:ellipse:rstar-compuesto}
\end{equation}
Interchanging the semiaxes transforms
\((\eta_{\mathrm{el}},R_\star,\tau_{\mathrm{el}})\) into
\((-\eta_{\mathrm{el}},R_\star^{-1},-1/\tau_{\mathrm{el}})\)
and preserves \(a_Sb_S=2\hbar\).
\end{proposition}
\begin{proof}
Dividing the semiaxes gives \(b_S/a_S=e^{-\eta_{\mathrm{el}}}\).
The map \(R\mapsto R^2\) is bijective on the positive
reals, so the positive root is unique.
For \(\xi=C^*/A\in(-1,1)\),
\[
 2\operatorname{artanh}\xi
 =\log\frac{1+\xi}{1-\xi};
\]
hence \(R_\star^4=(1-\xi)/(1+\xi)\).
Solving for \(\xi\) proves the first identity in
\eqref{x_reciprocidad:iv:ellipse:rstar-compuesto}. Substituting
\(A=1000\alpha\) and
\(C^*=2(\eta_{\mathrm{ret}}+\alpha)\) gives
\[
 \frac{A-C^*}{A+C^*}
 =\frac{998\alpha-2\eta_{\mathrm{ret}}}
        {1002\alpha+2\eta_{\mathrm{ret}}}
 =\frac{499\alpha-\eta_{\mathrm{ret}}}
        {501\alpha+\eta_{\mathrm{ret}}}.
\]
The interchange inverts the semiaxis ratio and changes the
sign of \(\eta_{\mathrm{el}}\). The identity
\(-1/(iR_\star^2)=iR_\star^{-2}\) gives the modular
transport. The product of the semiaxes remains fixed.
\end{proof}

On the positive branch, \(0<R_\star<1\); the chart of the
reversed form publishes \(R_\star^{-1}>1\). The angular ratio
is recovered from the radius, while the absolute angles
and their history remain in the enriched state.
The circular form corresponds to \(C^*/A=0\), with
\(R_\star=1\). Changing the positive action section
while retaining the angular pair changes the common scale of
the semiaxes and preserves their ratio.

\begin{proposition}[Lattice form determined by the ellipse]
\label{x_reciprocidad:iv:ellipse:matriz-prop}
The matrix of normalized squared semiaxes is
\begin{equation}
 D(\eta_{\mathrm{el}})
 =\frac1{2\hbar}
    \begin{pmatrix}a_S^2&0\\0&b_S^2\end{pmatrix}
 =\begin{pmatrix}R_\star^{-2}&0\\0&R_\star^2\end{pmatrix}.
 \label{x_reciprocidad:iv:ellipse:matriz}
\end{equation}
Its evaluation on an integer pair produces
\begin{equation}
 E_\star(m,w)=
 \begin{pmatrix}m&w\end{pmatrix}
 D(\eta_{\mathrm{el}})
 \begin{pmatrix}m\\w\end{pmatrix}
 =\frac{m^2}{R_\star^2}+w^2R_\star^2.
 \label{x_reciprocidad:iv:ellipse:energia}
\end{equation}
The equation of the ellipse multiplied by \(2\hbar\)
uses, instead, the inverse matrix \(D(-\eta_{\mathrm{el}})\).
\end{proposition}
\begin{proof}
The squares of the semiaxes in
\eqref{x_reciprocidad:iv:ellipse:semiejes}, divided by \(2\hbar\),
are \(e^{\eta_{\mathrm{el}}}\) and
\(e^{-\eta_{\mathrm{el}}}\).
By \eqref{x_reciprocidad:iv:ellipse:rstar}, these are
\(R_\star^{-2}\) and \(R_\star^2\).
Multiplication by the integer vector proves
\eqref{x_reciprocidad:iv:ellipse:energia}. Finally,
\(2\hbar/a_S^2=e^{-\eta_{\mathrm{el}}}\) and
\(2\hbar/b_S^2=e^{\eta_{\mathrm{el}}}\), identifying
the matrix of the elliptic equation.
\end{proof}

The radius \(R_\star\) is a dimensionless publication of
the labelled form. The polar radius
\(\rho_{\partial}\) belongs to the square-root-of-action line.
Representation by a spatial radius requires its subsequent
dimensional chart; this preserves the form already determined
and the independence of its construction from a
prescribed length.
